\chapter{Livšic theory}
\label{ch:livsic}

The rigidity theory of the preceding chapters concerned \emph{spectral} invariants --- Lyapunov
exponents, entropies, dimensions. This chapter develops a different kind of rigidity, the
\emph{cohomological} rigidity of A.~N.~Livšic (1972): over a hyperbolic-type map \(T\), a Hölder
observable \(\varphi\) is a coboundary --- \(\varphi = u\circ T - u\) for a Hölder transfer function
\(u\) --- \emph{if and only if} its Birkhoff sum vanishes around every periodic orbit. The forward
implication is a one-line telescoping identity; the converse reconstructs \(u\) from the periodic
data alone, and is the substance of the theorem. It is the exact obstruction theory for the
cohomological equation \(\varphi = u\circ T - u\): the countable family of closed-orbit sums is a
\emph{complete} set of invariants for solvability.

The library builds the full ladder. An \emph{abstract} Hölder Livšic theorem
(\Cref{thm:livsic-iff}) is proved from three regularity-free ingredients --- a dense forward orbit,
a Hölder observable, and a single geometric hypothesis, the \emph{exponential closing property}
(\Cref{def:livsic-expclosing}) --- following the dense-orbit construction of Katok--Hasselblatt
(\emph{Introduction to the Modern Theory of Dynamical Systems}, Theorem 19.2.1). That abstraction is
then discharged on five concrete systems: the one-sided full shift, the two-sided (invertible) full
shift, subshifts of finite type (with the golden-mean shift as headline), the doubling map of the
circle, and --- the culmination --- the Arnold cat map, a genuine Anosov diffeomorphism. On top of
the existence theorem sit the \emph{measurable rigidity} tiers, which weaken the transfer function
from Hölder to merely measurable: a continuous tier, a bounded-measurable tier, and the full
unbounded-measurable tier (Katok--Hasselblatt 19.2.4), the last routed through the two-sided natural
extension with stable/unstable essential-oscillation bounds. Finally a \emph{flow} tier ports the
obstruction direction to suspension flows, where it is exercised on the cat-map suspension.

Everything here is formalized sorry-free on the same substrate as the multiplicative theory and is
verified by the guarded axiom audit to rest only on
\(\{\texttt{propext},\texttt{Classical.choice},\texttt{Quot.sound}\}\). Throughout, \(T:X\to X\) is
a self-map of a metric space, \(\varphi:X\to\R\) an observable, and \(S_n\varphi\,x =
\sum_{k<n}\varphi(T^k x)\) its Birkhoff sum (Mathlib's \(\texttt{birkhoffSum}\)). This is, to our
knowledge, the first formalization of the Livšic theorem in any proof assistant.

\section{The abstract Hölder Livšic theorem}

\begin{definition}[Coboundary]\label{def:livsic-coboundary}
  \lean{ErgodicTheory.IsCoboundary}\leanok
  An observable \(\varphi\) is a \emph{coboundary} for a bare map \(T:X\to X\) if
  \(\varphi = u\circ T - u\) for some transfer function \(u:X\to\R\), with no regularity assumed on
  \(u\).
\end{definition}

\begin{definition}[Hölder coboundary]\label{def:livsic-holdercoboundary}
  \lean{ErgodicTheory.IsHolderCoboundary}\leanok
  \uses{def:livsic-coboundary}
  On a metric space, \(\varphi\) is a \emph{Hölder coboundary} for \(T\) if
  \(\varphi = u\circ T - u\) with \(u\) Hölder continuous (some exponent \(0 < r\), constant \(C\)).
  This is the target regularity of the Livšic theorem; forgetting the modulus of \(u\) shows a
  Hölder coboundary is a coboundary.
\end{definition}

\begin{definition}[Vanishing periodic sums]\label{def:livsic-vps}
  \lean{ErgodicTheory.HasVanishingPeriodicSums}\leanok
  \(\varphi\) \emph{has vanishing periodic sums} for \(T\) if every Birkhoff sum around a periodic
  orbit vanishes: whenever \(T^n p = p\), one has \(S_n\varphi\,p = \sum_{i<n}\varphi(T^i p) = 0\).
  This is the obstruction class the Livšic theorem certifies as complete.
\end{definition}

\begin{lemma}[Telescoping identity]\label{lem:livsic-telescope}
  \lean{ErgodicTheory.birkhoffSum_eq_of_coboundary}\leanok
  \uses{def:livsic-coboundary}
  If \(\varphi\,x = u(T x) - u\,x\) for all \(x\), then the Birkhoff sum collapses to the endpoints,
  \(S_n\varphi\,x = u(T^n x) - u\,x\). Pure algebra --- no metric, no regularity, valid for a bare
  \(T\).
\end{lemma}
\begin{proof}\leanok
  \uses{def:livsic-coboundary}
  Induction on \(n\) through \(\texttt{birkhoffSum\_succ}\): the successor term \(\varphi(T^k x) =
  u(T^{k+1}x) - u(T^k x)\) cancels against the running endpoint, leaving \(u(T^n x) - u\,x\).
\end{proof}

\begin{theorem}[Trivial direction and its obstruction certificate]\label{thm:livsic-trivial}
  \lean{ErgodicTheory.IsCoboundary.hasVanishingPeriodicSums}\leanok
  \uses{def:livsic-coboundary, def:livsic-vps, lem:livsic-telescope}
  A coboundary has vanishing periodic sums. Contrapositively, a single non-vanishing periodic sum
  defeats \emph{every} coboundary: if \(T^n p = p\) and \(S_n\varphi\,p\ne 0\), then \(\varphi\) is
  not a coboundary of any transfer function
  (\lean{ErgodicTheory.not_isCoboundary_of_periodicSum_ne_zero}).
\end{theorem}
\begin{proof}\leanok
  \uses{lem:livsic-telescope, def:livsic-vps}
  Apply the telescoping identity (\Cref{lem:livsic-telescope}) at a periodic point:
  \(S_n\varphi\,p = u(T^n p) - u\,p = u\,p - u\,p = 0\). The certificate is its contrapositive.
  Both hold for a bare \(T:X\to X\) and any \(u\); this is the downstream ``no continuous section''
  obstruction.
\end{proof}

The converse is the theorem. Its geometric engine is an abstracted closing property, stated in
\emph{summed-bound} form so that a single abstraction covers both the two-sided Anosov and the
one-sided/expanding shadowing regimes.

\begin{definition}[Exponential closing property]\label{def:livsic-expclosing}
  \lean{ErgodicTheory.ExpClosing}\leanok
  \(T\) has the \emph{exponential closing property} \(\texttt{ExpClosing}\,T\,\alpha\,\delta\,K\) if
  every point \(x\) that almost \(n\)-returns (\(\operatorname{dist}(x,T^n x)\le\delta\)) is shadowed by a genuine
  \(n\)-periodic point \(p\) (\(T^n p = p\)) whose orbit \(\alpha\)-Hölder-shadows that of \(x\) with
  \emph{total} cost controlled by the return gap:
  \[
    \sum_{i<n}\operatorname{dist}(T^i x, T^i p)^{\alpha} \;\le\; K\cdot\operatorname{dist}(x,T^n x)^{\alpha}.
  \]
  Summing the per-step cost (rather than hardcoding a per-term \(\min\)-exponent bound) is
  deliberate: the individual terms decay like \(\theta^{\min(i,\,n-i)}\) (two-sided Anosov) or
  \(\theta^{\,n-i}\) (one-sided/expanding), and both sum to a constant multiple of the endpoint gap.
  Abstracting the sum keeps the crux estimate regime-agnostic.
\end{definition}

\begin{lemma}[The crux estimate]\label{lem:livsic-crux}
  \lean{ErgodicTheory.abs_birkhoffSum_le}\leanok
  \uses{def:livsic-expclosing, def:livsic-vps}
  Let \(\varphi\) be \(r\)-Hölder with constant \(C_\varphi\), let \(T\) satisfy
  \(\texttt{ExpClosing}\,T\,r\,\delta\,K\), and let \(\varphi\) have vanishing periodic sums. Then
  any almost-\(n\)-return is controlled by its gap,
  \[
    \abs{S_n\varphi\,x} \;\le\; C_\varphi\,K\cdot\operatorname{dist}(x,T^n x)^{r},
    \qquad\text{whenever } \operatorname{dist}(x,T^n x)\le\delta.
  \]
\end{lemma}
\begin{proof}\leanok
  \uses{def:livsic-expclosing, def:livsic-vps}
  Subtract the (vanishing) periodic Birkhoff sum of the shadowing point \(p\) from that of \(x\):
  \(S_n\varphi\,x = \sum_{i<n}\bigl(\varphi(T^i x) - \varphi(T^i p)\bigr)\). Bound each term by
  Hölder continuity, \(\abs{\varphi(T^i x) - \varphi(T^i p)}\le C_\varphi\operatorname{dist}(T^i x,T^i p)^{r}\),
  and collapse the summed shadowing cost through the summed closing bound. The closing exponent and
  the Hölder exponent are the same \(r\), so the geometric bookkeeping matches term by term with no
  constraint on \(r\).
\end{proof}

\begin{lemma}[McShane--Whitney Hölder extension]\label{lem:livsic-mcshane}
  \lean{ErgodicTheory.exists_holderWith_extension}\leanok
  A real-valued function that is \(\texttt{HolderOnWith}\ C\ r\) (exponent \(r\le 1\)) on a subset
  \(s\) of a metric space extends to a global \(\texttt{HolderWith}\ C\ r\) function agreeing with it
  on \(s\), with the \emph{same} constant and exponent.
\end{lemma}
\begin{proof}\leanok
  The McShane infimal-convolution formula \(v(x) = \inf_{y\in s}\bigl(u(y) + C\operatorname{dist}(x,y)^{r}\bigr)\).
  For \(r\le 1\) the modulus \(t\mapsto C t^{r}\) is subadditive, so the infimum is finite, agrees
  with \(u\) on \(s\), and inherits the Hölder bound with the same constant.
\end{proof}

\begin{theorem}[Abstract existence, Katok--Hasselblatt 19.2.1]\label{thm:livsic-existence}
  \lean{ErgodicTheory.exists_holderCoboundary_of_denseOrbit}\leanok
  \uses{def:livsic-holdercoboundary, def:livsic-vps, def:livsic-expclosing, lem:livsic-crux, lem:livsic-mcshane}
  Let \(T\) be continuous on a compact metric space \(X\), let \(\varphi\) be \(r\)-Hölder
  (\(0 < r\le 1\)), suppose \(T\) has the summed exponential closing property, and suppose some
  point \(x_0\) has a dense forward orbit. If \(\varphi\) has vanishing periodic sums, then
  \(\varphi\) is a Hölder coboundary.
\end{theorem}
\begin{proof}\leanok
  \uses{lem:livsic-crux, lem:livsic-mcshane, thm:livsic-trivial}
  The classical dense-orbit construction. Enumerate the orbit \(e_n = T^n x_0\) and set the running
  transfer values \(w_n = S_n\varphi\,x_0\); define \(u_0\) on the orbit by \(u_0(e_m) = w_m\). This
  is well-defined because two indices landing on the same orbit point give equal Birkhoff sums (the
  difference is a Birkhoff sum around a periodic orbit, which vanishes). The close-pair increment
  \(\abs{w_b - w_a}\) is a Birkhoff sum over the segment \([a,b)\), so the crux estimate
  (\Cref{lem:livsic-crux}) gives it a Hölder modulus on the orbit; compactness bounds \(u_0\), which
  upgrades the modulus to a genuine \(\texttt{HolderOnWith}\). Extend to a global Hölder \(v\) by
  McShane (\Cref{lem:livsic-mcshane}). The on-orbit cohomology identity \(\varphi(e_m) = v(e_{m+1}) -
  v(e_m)\) (with \(e_{m+1} = T e_m\)) then transports to all of \(X\) by density and continuity
  (\(\texttt{Continuous.ext\_on}\)), giving \(\varphi = v\circ T - v\).
\end{proof}

\begin{theorem}[The Livšic equivalence]\label{thm:livsic-iff}
  \lean{ErgodicTheory.isHolderCoboundary_iff}\leanok
  \uses{thm:livsic-trivial, thm:livsic-existence}
  Under the standing hypotheses (continuous \(T\) on a compact metric space, \(r\)-Hölder
  \(\varphi\) with \(0 < r\le 1\), summed exponential closing, and a dense forward orbit),
  \[
    \varphi \text{ is a Hölder coboundary} \iff \varphi \text{ has vanishing periodic sums}.
  \]
\end{theorem}
\begin{proof}\leanok
  \uses{thm:livsic-trivial, thm:livsic-existence}
  Forward: the pure telescoping obstruction (\Cref{thm:livsic-trivial}). Backward: the dense-orbit
  reconstruction (\Cref{thm:livsic-existence}). This is the interface every concrete instance below
  discharges by supplying continuity, compactness, a closing constant, and a dense orbit.
\end{proof}

\section{Concrete instances}

Each instance of \Cref{thm:livsic-iff} needs a metric substrate, continuity, compactness, an
\(\texttt{ExpClosing}\) constant, and a dense forward orbit. The instances below range from the
combinatorially trivial full shift to a genuine Anosov diffeomorphism.

\subsection{The one-sided full shift}

The full shift \(\sigma\) on \(\texttt{Shift}\ \alpha_0 = (\N\to\alpha_0)\) over a finite alphabet
carries the \(\texttt{PiNat}\) ultrametric with \(\operatorname{dist}(x,y) = (1/2)^{N}\) at first disagreement
index \(N\). It carries \emph{no admissibility bookkeeping} --- every finite word is legal --- so
the closing property is available unconditionally.

\begin{lemma}[Shift metric substrate]\label{lem:livsic-shiftmetric}
  \lean{ErgodicTheory.Livsic.lipschitzWith_two_shiftMap}\leanok
  The left shift is \(2\)-Lipschitz for the \(\texttt{PiNat}\) ultrametric
  (\lean{ErgodicTheory.Livsic.instIsUltrametricDist_shift}), hence continuous; the full shift over a
  finite discrete alphabet is compact by Tychonoff.
\end{lemma}
\begin{proof}\leanok
  Dropping the first coordinate moves the first-disagreement index down by one, doubling the
  distance at worst; compactness is a product of finite discrete spaces.
\end{proof}

\begin{theorem}[Closing for the full shift]\label{thm:livsic-shift-closing}
  \lean{ErgodicTheory.Livsic.expClosing_shiftMap}\leanok
  \uses{def:livsic-expclosing}
  For every exponent \(\alpha > 0\) the full shift satisfies \(\texttt{ExpClosing}\,\sigma\,\alpha\,1\,K\)
  with \(K = (1/2)^{\alpha}/(1 - (1/2)^{\alpha})\).
\end{theorem}
\begin{proof}\leanok
  \uses{def:livsic-expclosing}
  A point that almost \(n\)-returns is shadowed by the front-anchored periodization
  \(p_i := x_{i\bmod n}\), automatically admissible on the full shift. The \(i\)-th shifted pair
  agrees on the first \(n-i\) coordinates, so \(\operatorname{dist}(\sigma^i x,\sigma^i p)\le(1/2)^{\,n-i}\); the
  \(\alpha\)-powers form a geometric series summing to \(K\cdot\operatorname{dist}(x,\sigma^n x)^{\alpha}\).
\end{proof}

\begin{definition}[The rich point]\label{def:livsic-richpoint}
  \lean{ErgodicTheory.Livsic.richPoint}\leanok
  The \emph{rich point} is the sequence obtained by concatenating all finite words (enumerated
  through the alphabet's countable encoding). Every finite word appears as a block, so every target
  is approximated by some forward iterate.
\end{definition}

\begin{theorem}[Dense orbit for the full shift]\label{thm:livsic-shift-dense}
  \lean{ErgodicTheory.Livsic.exists_denseRange_shiftMap_orbit_alphabet}\leanok
  \uses{def:livsic-richpoint}
  Over a nonempty finite alphabet the forward orbit of the rich point is dense.
\end{theorem}
\begin{proof}\leanok
  \uses{def:livsic-richpoint}
  To match a target on its first \(N\) coordinates, shift the rich point to the block equal to that
  length-\(N\) prefix; the resulting iterate lies within \((1/2)^{N}\) of the target.
\end{proof}

\begin{theorem}[Livšic for the one-sided full shift]\label{thm:livsic-fullshift}
  \lean{ErgodicTheory.Livsic.livsic_fullShift}\leanok
  \uses{thm:livsic-iff, lem:livsic-shiftmetric, thm:livsic-shift-closing, thm:livsic-shift-dense}
  For a Hölder \(\varphi\) (exponent \(0 < r\le 1\)) on the full shift over a finite discrete
  alphabet,
  \[
    \varphi \text{ is a Hölder coboundary for } \sigma \iff \varphi \text{ has vanishing periodic sums}.
  \]
\end{theorem}
\begin{proof}\leanok
  \uses{thm:livsic-iff, lem:livsic-shiftmetric, thm:livsic-shift-closing, thm:livsic-shift-dense}
  Feed \Cref{thm:livsic-iff} the substrate (\Cref{lem:livsic-shiftmetric}), the closing constant
  (\Cref{thm:livsic-shift-closing}, \(\delta = 1\)), and the dense orbit
  (\Cref{thm:livsic-shift-dense}).
\end{proof}

The instance is non-vacuous on the binary shift. The locally constant potential
\(\phi\,x = \mathbf 1[x_0 = 1]\) is \emph{not} a coboundary
(\lean{ErgodicTheory.Livsic.phi_not_isCoboundary}): its period-\(1\) sum at the fixed point
\(\underline 1\) is \(1\ne 0\), so the obstruction certificate of \Cref{thm:livsic-trivial} applies.
By contrast \(\psi = \phi\circ\sigma - \phi\) \emph{is} a Hölder coboundary
(\lean{ErgodicTheory.Livsic.psi_isHolderCoboundary}) by construction, and re-deriving that fact
\emph{through} the hard backward direction of \Cref{thm:livsic-fullshift} exercises the dense-orbit
reconstruction end-to-end on a concrete witness
(\lean{ErgodicTheory.Livsic.psi_isHolderCoboundary_via_livsic}).

\subsection{The two-sided full shift}

The invertible finale requires a bi-infinite metric. Rather than transport the \(\texttt{PiNat}\)
metric through a \(\Z\simeq\N\) reindexing (which scrambles the cylinder dictionary), the library
builds the \(\Z\)-indexed \(\theta\)-ultrametric \emph{directly}, measuring first disagreement by
absolute value of the integer coordinate.

\begin{definition}[The \(\Z\)-indexed \(\theta\)-ultrametric]\label{def:livsic-bishift-metric}
  \lean{ErgodicTheory.biShiftMetricSpace}\leanok
  On \(\texttt{BiShift}\ \alpha_0 = (\Z\to\alpha_0)\), set \(\operatorname{dist}(x,y) = (1/2)^{N}\) where
  \(N\) is the least \(\abs{j}\) at which \(x_j\ne y_j\) (\lean{ErgodicTheory.distZ}); the symmetric
  cylinders \(\{y : \forall\,\abs{j}<N,\ y_j = x_j\}\) are the balls. The metric is built against
  the pre-existing product topology, so \(\texttt{CompactSpace}\) and \(\texttt{BorelSpace}\) are
  inherited with no diamond.
\end{definition}

\begin{lemma}[Two-sided geometric bound]\label{lem:livsic-minregime}
  \lean{ErgodicTheory.min_regime_sum_le}\leanok
  For \(0\le\theta<1\), \(\ \sum_{i<n}\theta^{\min(i,\,n-i)} \le 2\,(1-\theta)^{-1}\).
\end{lemma}
\begin{proof}\leanok
  The two-sided profile \(\theta^{\min(i,\,n-i)}\) is dominated termwise by \(\theta^{i} +
  \theta^{\,n-i}\); the first arm is a geometric partial sum and the second reflects to another one,
  giving \emph{twice} the one-sided bound.
\end{proof}

\begin{theorem}[Closing for the two-sided full shift]\label{thm:livsic-bishift-closing}
  \lean{ErgodicTheory.expClosing_biShiftMap}\leanok
  \uses{def:livsic-expclosing, def:livsic-bishift-metric, lem:livsic-minregime}
  For every \(\alpha>0\) the two-sided full shift \(\tilde\sigma\) is \(2\)-Lipschitz
  (\lean{ErgodicTheory.lipschitzWith_two_biShiftMap}) and satisfies
  \(\texttt{ExpClosing}\,\tilde\sigma\,\alpha\,1\,K\) with \(K = 2/(1 - (1/2)^{\alpha})\) --- twice
  the one-sided constant.
\end{theorem}
\begin{proof}\leanok
  \uses{def:livsic-expclosing, lem:livsic-minregime}
  The shadow of an almost-\(n\)-return is the \emph{central} periodization \(p_j := x_{j\bmod n}\)
  (integer \(\bmod\)), which repeats the central block \(x_0,\dots,x_{n-1}\). If \(x\) and
  \(\tilde\sigma^n x\) first differ at \(\abs{\cdot}=N\), then \(x\) is genuinely periodic on the
  whole window \(-N < j < n+N\), so \(p\) agrees with \(x\) there and the per-step shadowing distance
  follows the bilateral profile \(\theta^{\min(i,\,n-i)}\); \Cref{lem:livsic-minregime} sums it.
\end{proof}

\begin{theorem}[Livšic for the two-sided full shift]\label{thm:livsic-bishift}
  \lean{ErgodicTheory.livsic_biShift}\leanok
  \uses{thm:livsic-iff, def:livsic-bishift-metric, thm:livsic-bishift-closing}
  For a Hölder \(\varphi\) (\(0 < r\le 1\)) on the two-sided full shift, \(\varphi\) is a Hölder
  coboundary for the bilateral shift iff all its periodic Birkhoff sums vanish. The dense forward
  orbit is \lean{ErgodicTheory.exists_denseRange_biShiftMap_orbit} (the one-sided rich point placed
  on the non-negative axis, padded on the negative axis).
\end{theorem}
\begin{proof}\leanok
  \uses{thm:livsic-iff, thm:livsic-bishift-closing}
  Instantiate \Cref{thm:livsic-iff} with the \(\Z\)-ultrametric, the doubled closing constant, and
  the two-sided dense orbit. Compactness comes with no metric/measure diamond because the metric
  reuses the product topology whose Borel structure carries the two-sided Bernoulli measure.
\end{proof}

\subsection{Subshifts of finite type and the golden-mean shift}

A subshift of finite type (SFT) is cut out of the full shift by a \(0/1\) transition matrix
\(M\). Here the closing property becomes delicate --- the front-anchored periodization has a wrap
transition that need not be legal --- and the library isolates a clean design finding that dissolves
the difficulty.

\begin{definition}[The SFT carrier and shift]\label{def:livsic-sft}
  \lean{ErgodicTheory.SFTCarrier}\leanok
  For \(M:\Fin k\to\Fin k\to\texttt{Bool}\), the \emph{carrier}
  \(\texttt{SFTCarrier}\ M = \{x : \forall i,\ M(x_i, x_{i+1})\}\) is the closed shift-invariant set
  of admissible sequences; as a subtype (\lean{ErgodicTheory.SFT}) it inherits the ultrametric and,
  being closed in the compact full shift, is compact. The \emph{SFT shift}
  \lean{ErgodicTheory.sftShiftMap} is the restriction of \(\sigma\), still \(2\)-Lipschitz.
\end{definition}

\begin{lemma}[The \(\delta = 1/2\) crux]\label{lem:livsic-sft-half}
  \lean{ErgodicTheory.sft_head_eq_of_dist_le_half}\leanok
  \uses{def:livsic-sft}
  If \(x\) almost \(n\)-returns within radius \(1/2\), then \(x_0 = x_n\).
\end{lemma}
\begin{proof}\leanok
  \uses{def:livsic-sft}
  \(\operatorname{dist}(x,\sigma^n x)\le 1/2\) forces agreement on coordinate \(0\), i.e.\ \(x_0 = (\sigma^n x)_0 =
  x_n\).
\end{proof}

\begin{theorem}[Closing for the SFT, unconditional in \(M\)]\label{thm:livsic-sft-closing}
  \lean{ErgodicTheory.expClosing_sftShiftMap}\leanok
  \uses{def:livsic-expclosing, def:livsic-sft, lem:livsic-sft-half}
  For every \(\alpha>0\) and \emph{every} transition matrix \(M\), the SFT shift satisfies
  \(\texttt{ExpClosing}\,(\texttt{sftShiftMap}\ M)\,\alpha\,(1/2)\,K\) at the reduced closing radius
  \(\delta = 1/2\), with the same constant \(K = (1/2)^{\alpha}/(1-(1/2)^{\alpha})\) as the full
  shift --- \emph{with no irreducibility hypothesis on \(M\)}.
\end{theorem}
\begin{proof}\leanok
  \uses{lem:livsic-sft-half}
  The design finding: at radius \(1/2\) the periodization \(p_i := x_{i\bmod n}\) is admissible for
  \emph{any} \(M\). By \Cref{lem:livsic-sft-half} the wrap transition \(M(x_{n-1}, x_0) =
  M(x_{n-1}, x_n)\) coincides with the interior transition \(M(x_{n-1}, x_{(n-1)+1})\), which is
  legal because \(x\) itself is admissible. No connecting word is needed. The shadow bound then
  transports the full-shift estimate through the subtype metric.
\end{proof}

Irreducibility is used \emph{only} to build a dense orbit. The library isolates exactly the
combinatorial input the dense-orbit construction needs, disclosed honestly as strictly stronger than
irreducibility.

\begin{definition}[Safe symbol]\label{def:livsic-safesymbol}
  \lean{ErgodicTheory.SafeSymbol}\leanok
  A \emph{safe symbol} for \(M\) is a symbol \(s\) allowed adjacent to every symbol:
  \(M(a,s)\) and \(M(s,a)\) hold for all \(a\). This is \emph{strictly stronger} than irreducibility
  of \(M\); a general irreducible SFT need not have one, and dense orbits for such SFTs are
  \textbf{not} delivered here. The golden-mean shift has \(s = 0\).
\end{definition}

\begin{theorem}[Dense orbit for an SFT with a safe symbol]\label{thm:livsic-sft-dense}
  \lean{ErgodicTheory.exists_denseRange_sftShiftMap_orbit}\leanok
  \uses{def:livsic-safesymbol, def:livsic-sft}
  If \(M\) has a safe symbol, the SFT shift has a point with dense forward orbit.
\end{theorem}
\begin{proof}\leanok
  \uses{def:livsic-safesymbol}
  Mirror the full-shift rich point, but first \emph{sanitize} each decoded word (keep it if
  admissible, drop it otherwise) and pad every block with one trailing safe symbol. The pad makes
  every inter-block junction legal, so the whole sequence lands in the carrier; every admissible
  target word is retained verbatim as its own block, giving density.
\end{proof}

\begin{theorem}[Livšic for a one-sided SFT]\label{thm:livsic-sft}
  \lean{ErgodicTheory.livsic_sft}\leanok
  \uses{thm:livsic-iff, def:livsic-sft, thm:livsic-sft-closing}
  For a Hölder \(\varphi\) (\(0 < r\le 1\)) on the SFT and a dense forward orbit under the SFT shift,
  \(\varphi\) is a Hölder coboundary iff all its periodic Birkhoff sums vanish.
\end{theorem}
\begin{proof}\leanok
  \uses{thm:livsic-iff, thm:livsic-sft-closing}
  Instantiate \Cref{thm:livsic-iff} with the subtype substrate and the unconditional
  \(\delta = 1/2\) closing (\Cref{thm:livsic-sft-closing}); the dense orbit is the sole hypothesis,
  the only place a combinatorial assumption on \(M\) enters.
\end{proof}

\begin{theorem}[The golden-mean shift, unconditional]\label{thm:livsic-goldenmean}
  \lean{ErgodicTheory.livsic_goldenMean}\leanok
  \uses{thm:livsic-sft, thm:livsic-sft-dense, def:livsic-safesymbol}
  For the golden-mean shift (\lean{ErgodicTheory.goldenMeanM}: forbid the block \(11\), Lind--Marcus
  \S1.2), a Hölder \(\varphi\) (\(0 < r\le 1\)) is a Hölder coboundary iff all its periodic Birkhoff
  sums vanish --- with no external hypotheses.
\end{theorem}
\begin{proof}\leanok
  \uses{thm:livsic-sft, thm:livsic-sft-dense}
  The symbol \(0\) is safe (it forbids nothing), so \Cref{thm:livsic-sft-dense} supplies the dense
  orbit and \Cref{thm:livsic-sft} closes the equivalence. The shift is a \emph{proper} subshift ---
  the all-ones sequence contains the forbidden block, \lean{ErgodicTheory.goldenMean_proper} --- so
  the instance is not the full shift in disguise.
\end{proof}

\subsection{The doubling map}

The doubling map \(y\mapsto 2y\) on the circle is the canonical smooth expanding endomorphism. Its
dense orbit comes from a generic ergodicity lemma, and its closing is an explicit rounding
construction.

\begin{theorem}[Dense orbit from ergodicity]\label{thm:livsic-ergodic-dense}
  \lean{ErgodicTheory.ergodic_exists_denseRange_iterate}\leanok
  On a second-countable space with an open-positive probability measure, an ergodic map has a point
  with dense forward orbit.
\end{theorem}
\begin{proof}\leanok
  For each basic open \(o\) the visiting set \(U_o = \bigcup_n T^{-n}o\) is almost invariant, hence
  null or conull by ergodicity; since \(o\subseteq U_o\) has positive measure it is conull.
  Intersecting over the countable basis leaves a conull set of points whose orbit visits every basic
  open set, i.e.\ a dense orbit; nonemptiness follows from the probability normalization. This
  generic lemma is Mathlib-upstreamable and is reused for the cat map.
\end{proof}

\begin{theorem}[Closing for the doubling map]\label{thm:livsic-doubling-closing}
  \lean{ErgodicTheory.expClosing_doublingMap}\leanok
  \uses{def:livsic-expclosing}
  For every \(r>0\) the doubling map satisfies \(\texttt{ExpClosing}\ r\ 1\ (2^{r}/(2^{r}-1))\).
\end{theorem}
\begin{proof}\leanok
  \uses{def:livsic-expclosing}
  Because \(T^n = (2^n)\cdot\) is a group endomorphism, the periodic shadow of \(x = \bar a\) is
  obtained by \emph{rounding}: with \(M = 2^n - 1\) and \(k = \operatorname{round}(Ma)\), the point
  \(p = \bar a - \overline{(Ma-k)/M}\) is genuinely \(n\)-periodic (\(M\cdot p = 0\)) and lies within
  \(\norm{x - 2^n x}/M\) of \(x\) (\lean{ErgodicTheory.exists_doubling_periodic_shadow}). The
  \(i\)-th orbit points satisfy \(\operatorname{dist}(2^i x, 2^i p)\le 2^i\norm{x-p}\), and the \(r\)-powers sum to
  a geometric series bounded via \(2^n/(2^n-1)\le 2\).
\end{proof}

\begin{theorem}[Livšic for the doubling map]\label{thm:livsic-doubling}
  \lean{ErgodicTheory.livsic_doublingMap}\leanok
  \uses{thm:livsic-iff, thm:livsic-ergodic-dense, thm:livsic-doubling-closing}
  For a Hölder \(\varphi\) (\(0 < r\le 1\)) on the circle, \(\varphi\) is a Hölder coboundary for the
  doubling map iff all its periodic Birkhoff sums vanish.
\end{theorem}
\begin{proof}\leanok
  \uses{thm:livsic-iff, thm:livsic-ergodic-dense, thm:livsic-doubling-closing}
  The doubling map is continuous, the circle compact, the closing is
  \Cref{thm:livsic-doubling-closing}, and the dense orbit is \Cref{thm:livsic-ergodic-dense} applied
  to the ergodicity of the doubling map for Haar measure
  (\lean{ErgodicTheory.exists_denseRange_doublingMap_orbit}). Non-vacuity: the constant \(1\) is not
  a coboundary (\lean{ErgodicTheory.const_one_not_isCoboundary_doublingMap}), its sum at the fixed
  point \(0\) being \(1\ne 0\).
\end{proof}

\subsection{The Arnold cat map}

The culmination is Livšic on a genuine Anosov diffeomorphism: the hyperbolic toral automorphism
\(\texttt{catTorus}:\mathbb{T}^2\to\mathbb{T}^2\) induced by \(\bigl(\begin{smallmatrix}2&1\\1&1\end{smallmatrix}\bigr)\).
Here the Anosov closing lemma admits an \emph{exact}, non-iterative solution.

\begin{definition}[Covering projection]\label{def:livsic-catproj}
  \lean{ErgodicTheory.CatMapToral.catProj}\leanok
  \(\texttt{catProj}:\R^2\to\mathbb{T}^2\) is the coordinatewise quotient projection; it intertwines the
  linear map \(\texttt{cat}_{\R} = \bigl(\begin{smallmatrix}2&1\\1&1\end{smallmatrix}\bigr)\)
  upstairs with \(\texttt{catTorus}\) downstairs and is distance-nonincreasing.
\end{definition}

\begin{definition}[Exact shadow solution]\label{def:livsic-catshadow}
  \lean{ErgodicTheory.CatMapToral.catShadowSol}\leanok
  For a lifted return defect \(e\), \(\texttt{catShadowSol}\ n\ e\) solves the cohomological
  equation \((\texttt{cat}_{\R}^n - 1)\,w = e\) in closed form in the hyperbolic eigenbasis --- \(1\)
  is not an eigenvalue of any power of the hyperbolic matrix, so \(\texttt{cat}_{\R}^n - 1\) is
  invertible on \(\R^2\).
\end{definition}

\begin{theorem}[Anosov exact-solve closing]\label{thm:livsic-cat-closing}
  \lean{ErgodicTheory.CatMapToral.expClosing_catTorus}\leanok
  \uses{def:livsic-expclosing, def:livsic-catproj, def:livsic-catshadow}
  For every exponent \(0 < \alpha\le 1\), \(\texttt{catTorus}\) has the summed exponential closing
  property with radius \(1\) and an explicit constant \(2\,C_{\mathrm{sh}}^{\alpha}/(1-\theta^{\alpha})\),
  where \(\theta<1\) is the contraction ratio (the small eigenvalue).
\end{theorem}
\begin{proof}\leanok
  \uses{def:livsic-catproj, def:livsic-catshadow}
  Lift an almost-return to \(d = \texttt{cat}_{\R}^n v - v\), reduce it to its nearest-integer
  representative \(e\), solve \(w = \texttt{catShadowSol}\ n\ e\), and project
  \(p = \texttt{catProj}(v - w)\). Then \(p\) is \emph{exactly} \(n\)-periodic because \(d - e\) is an
  integer vector (invisible to the projection), and its forward orbit two-sidedly geometrically
  shadows that of \(x\): the shadow displacement at step \(i\) is \(\norm{\texttt{cat}_{\R}^i w}\),
  whose eigenbasis decomposition gives the two-sided \(\theta^{\min(i,\,n-i)}\) profile with summed
  \(\alpha\)-power bounded by the return-gap constant.
\end{proof}

\begin{theorem}[Livšic for the Arnold cat map]\label{thm:livsic-cat}
  \lean{ErgodicTheory.CatMapToral.livsic_catTorus}\leanok
  \uses{thm:livsic-iff, thm:livsic-ergodic-dense, thm:livsic-cat-closing}
  For a Hölder \(\varphi\) (\(0 < r\le 1\)) on \(\mathbb{T}^2\), \(\varphi\) is a Hölder coboundary for the
  Anosov diffeomorphism \(\texttt{catTorus}\) iff all its periodic Birkhoff sums vanish.
\end{theorem}
\begin{proof}\leanok
  \uses{thm:livsic-iff, thm:livsic-ergodic-dense, thm:livsic-cat-closing}
  \(\texttt{catTorus}\) is continuous, \(\mathbb{T}^2\) compact, the closing is
  \Cref{thm:livsic-cat-closing}, and the dense orbit is \Cref{thm:livsic-ergodic-dense} applied to
  the genuine ergodicity of the cat map for Haar measure. Non-vacuity: the constant \(1\) is not a
  coboundary (\lean{ErgodicTheory.CatMapToral.const_one_not_isCoboundary_catTorus}), its sum at the
  fixed point \(0\) being \(1\ne 0\).
\end{proof}

\section{Measurable rigidity tiers}

The existence theorem produces a \emph{Hölder} transfer function. The measurable rigidity programme
asks the converse robustness question: if the cohomological equation is solved only
\emph{almost everywhere} by a merely \emph{measurable} \(u\), must \(\varphi\) already have vanishing
periodic sums (hence be a genuine Hölder coboundary)? This splits into three tiers by the regularity
assumed of \(u\).

\begin{definition}[Almost-everywhere coboundary]\label{def:livsic-aecob}
  \lean{ErgodicTheory.IsAeCoboundaryOf}\leanok
  \(\varphi\) is an \emph{a.e.\ coboundary} of \(u\) over \(\mu\) if \(\varphi = u\circ T - u\) holds
  \(\mu\)-almost-everywhere. This is the measurable analogue of \Cref{def:livsic-coboundary}.
\end{definition}

\subsection{Continuous and bounded tiers}

\begin{theorem}[Continuous tier]\label{thm:livsic-continuous-tier}
  \lean{ErgodicTheory.Livsic.hasVanishingPeriodicSums_of_continuous_coboundary}\leanok
  \uses{def:livsic-aecob, thm:livsic-trivial}
  If \(\varphi\) is continuous and equals the coboundary of a \emph{continuous} \(u\) only
  \(\mu\)-a.e., for a fully supported \(\mu\) (\(\texttt{IsOpenPosMeasure}\)), then \(\varphi\) has
  vanishing periodic sums.
\end{theorem}
\begin{proof}\leanok
  \uses{thm:livsic-trivial}
  Full support forces a.e.-equal continuous functions to be equal, so the a.e.\ equation upgrades to
  an \emph{everywhere} coboundary and \Cref{thm:livsic-trivial} finishes. No ergodicity or
  hyperbolicity is used. On the full shift, a fully supported Bernoulli law is open-positive, giving
  the corollary \lean{ErgodicTheory.Livsic.isHolderCoboundary_of_continuous_aeCoboundary}.
\end{proof}

\begin{theorem}[Bounded tier]\label{thm:livsic-bounded-tier}
  \lean{ErgodicTheory.Livsic.vanishingPeriodicSum_of_bounded_shadowing}\leanok
  \uses{def:livsic-aecob, def:livsic-vps}
  Let \(T\) preserve a probability measure \(\mu\), let \(p\) be \(n\)-periodic, and let
  \(\varphi = u\circ T - u\) \(\mu\)-a.e.\ with \(u\) \emph{bounded} (\(\abs u\le M\)). If each
  cylinder \(D_m\) around \(p\) has positive measure and shadows the Birkhoff sum of \(p\) up to a
  uniform additive constant \(B\), then \(S_n\varphi\,p = 0\).
\end{theorem}
\begin{proof}\leanok
  \uses{def:livsic-vps}
  Around the periodic orbit \(S_{nm}\varphi\,p = m\cdot S_n\varphi\,p\). A.e.\ telescoping gives
  \(S_{nm}\varphi\,x = u(T^{nm}x) - u\,x\), bounded by \(2M\) \emph{uniformly in \(m\)}. Picking a
  witness in the positive-measure \(D_m\) yields \(\abs{m\cdot S_n\varphi\,p}\le B + 2M\), a bound
  independent of \(m\); if \(S_n\varphi\,p\ne 0\) the left side is unbounded, a contradiction. The
  uniformity of \(B\) and \(2M\) in \(m\) is exactly where unboundedness of \(u\) would break the
  argument.
\end{proof}

Instantiated on the full shift with a fully supported Bernoulli law --- cylinders are charged, and a
Hölder observable shadows along a cylinder with a depth-independent constant --- the bounded core
gives \lean{ErgodicTheory.Livsic.hasVanishingPeriodicSums_of_bounded_aeCoboundary}, which promotes to
the Hölder-coboundary conclusion \lean{ErgodicTheory.Livsic.isHolderCoboundary_of_bounded_aeCoboundary}
through \Cref{thm:livsic-fullshift}.

\subsection{The full measurable tier via the natural extension}

For a genuinely unbounded measurable \(u\) the uniform endpoint control of the bounded tier breaks,
and the theorem becomes the classical Livšic \emph{regularity} theorem (Katok--Hasselblatt 19.2.4).
The library discharges it through the two-sided natural extension, and does so with two disclosed
proof innovations: there is \emph{no Hölder-version construction} of the transfer function and
\emph{no Birkhoff ergodic theorem} --- the oscillation bounds are obtained by reverse Fatou.

\begin{lemma}[The natural-extension factor]\label{lem:livsic-toshift}
  \lean{ErgodicTheory.toShift}\leanok
  The restriction \(\texttt{toShift}:\texttt{BiShift}\ \alpha_0\to\texttt{Shift}\ \alpha_0\) to the
  non-negative coordinates is a measure-preserving factor map
  (\lean{ErgodicTheory.measurePreserving_toShift}) intertwining the two-sided shift with the
  one-sided one: the one-sided Bernoulli system is a factor of the invertible two-sided one. Hölder
  data, measurability, and the a.e.\ cohomological equation all transport \emph{up} this factor.
\end{lemma}
\begin{proof}\leanok
  The factor map is \(1\)-Lipschitz for the respective \(\theta\)-ultrametrics and pushes the
  two-sided Bernoulli measure to the one-sided one; pulling the a.e.\ identity back along it and
  rewriting through the intertwining transports the coboundary equation.
\end{proof}

\begin{lemma}[Product structure]\label{lem:livsic-joinpf}
  \lean{ErgodicTheory.joinPF}\leanok
  The bilateral shift space splits measurably as \(\text{past}\otimes\text{future}\) via
  \(\texttt{joinPF}\), under which the two-sided Bernoulli measure is the product of the past and
  future Bernoulli laws (measure-preservingly).
\end{lemma}
\begin{proof}\leanok
  The past (\(j<0\)) and future (\(j\ge 0\)) coordinate blocks are independent under the i.i.d.\
  Bernoulli measure, so the joining map is measure preserving; this is the substrate for the Fubini
  glue below.
\end{proof}

\begin{lemma}[Stable and unstable essential oscillation]\label{lem:livsic-osc}
  \lean{ErgodicTheory.stable_pair_osc}\leanok
  \uses{def:livsic-aecob, lem:livsic-joinpf}
  For a measurable a.e.\ transfer function of an \(r\)-Hölder \(\varphi\) over the two-sided
  Bernoulli measure, and for a.e.\ \emph{stable} pair (two pasts, one shared future) the values of
  \(u\) differ by at most \(C\theta/(1-\theta)\); symmetrically for a.e.\ \emph{unstable} pair (two
  futures, one shared past), \lean{ErgodicTheory.unstable_pair_osc}.
\end{lemma}
\begin{proof}\leanok
  \uses{lem:livsic-joinpf}
  Points sharing a future have forward orbits that converge, so the telescoped Birkhoff sums of
  \(\varphi\) along them differ by a convergent geometric tail; Lusin regularity plus \emph{reverse
  Fatou} on the return set (in place of a Birkhoff density of returns) turns this into an a.e.\
  essential bound on \(u\)'s oscillation. The unstable bound is the \(\tilde\sigma^{-1}\) mirror,
  reversing the cocycle to \(\psi = -\varphi\circ\tilde\sigma^{-1}\).
\end{proof}

\begin{theorem}[Essential boundedness of the transfer]\label{thm:livsic-essbdd}
  \lean{ErgodicTheory.essBounded_of_measurable_aeCoboundary_biShift}\leanok
  \uses{lem:livsic-osc, lem:livsic-joinpf}
  A merely measurable (possibly unbounded) a.e.\ solution \(u\) of the cohomological equation over
  the two-sided Bernoulli measure is essentially bounded.
\end{theorem}
\begin{proof}\leanok
  \uses{lem:livsic-osc, lem:livsic-joinpf}
  The Fubini glue: over the \(\text{past}\otimes\text{future}\) product a typical base point
  \((a_0,b_0)\) connects to a.e.\ \((a,b)\) in two holonomy moves,
  \((a,b)\xrightarrow{\text{unstable}}(a,b_0)\xrightarrow{\text{stable}}(a_0,b_0)\), so
  \(\abs u\le\abs{u(a_0,b_0)} + C_s + C_u\) almost everywhere by the two oscillation bounds of
  \Cref{lem:livsic-osc}.
\end{proof}

\begin{theorem}[Two-sided measurable headline]\label{thm:livsic-measurable-bishift}
  \lean{ErgodicTheory.hasVanishingPeriodicSums_of_measurable_aeCoboundary_biShift}\leanok
  \uses{thm:livsic-essbdd, thm:livsic-bounded-tier}
  Over a fully supported two-sided Bernoulli measure, an \(r\)-Hölder \(\varphi\) that is the a.e.\
  coboundary of a merely measurable (possibly unbounded) \(u\) has vanishing periodic sums.
\end{theorem}
\begin{proof}\leanok
  \uses{thm:livsic-essbdd, thm:livsic-bounded-tier}
  With the essential bound \(M\) from \Cref{thm:livsic-essbdd}, the clamp \(v = \max(-M',\min(M',u))\)
  is a bounded-everywhere measurable function agreeing with \(u\) a.e., hence itself an a.e.\ transfer
  function (the shift being measure preserving); the already-shipped bounded tier
  (\Cref{thm:livsic-bounded-tier}, two-sided form) discharges vanishing periodic sums. No
  Hölder-version construction, no ergodic theorem.
\end{proof}

\begin{theorem}[One-sided measurable headline]\label{thm:livsic-measurable-oneside}
  \lean{ErgodicTheory.Livsic.hasVanishingPeriodicSums_of_measurable_aeCoboundary}\leanok
  \uses{thm:livsic-measurable-bishift, lem:livsic-toshift}
  Over a fully supported one-sided Bernoulli measure, an \(r\)-Hölder \(\varphi\) that is the a.e.\
  coboundary of a merely measurable transfer function has vanishing periodic sums.
\end{theorem}
\begin{proof}\leanok
  \uses{thm:livsic-measurable-bishift, lem:livsic-toshift}
  Transport \(\varphi,u\) up the natural-extension factor (\Cref{lem:livsic-toshift}); apply the
  two-sided headline (\Cref{thm:livsic-measurable-bishift}) to \(\varphi\circ\texttt{toShift}\); then
  \emph{descend} through a periodic lift --- a one-sided \(n\)-periodic point lifts to the
  bi-infinite \(n\)-periodization, whose lifted Birkhoff sum equals the one-sided sum, so vanishing
  descends.
\end{proof}

\begin{theorem}[Full measurable Livšic rigidity, Katok--Hasselblatt 19.2.4]\label{thm:livsic-measurable-full}
  \lean{ErgodicTheory.Livsic.livsic_measurable_rigidity}\leanok
  \uses{thm:livsic-measurable-oneside, thm:livsic-fullshift}
  Over a fully supported Bernoulli measure, an \(r\)-Hölder \(\varphi\) (\(0 < r\le 1\)) admits some
  \emph{measurable} a.e.\ transfer function \emph{iff} all its periodic Birkhoff sums vanish.
\end{theorem}
\begin{proof}\leanok
  \uses{thm:livsic-measurable-oneside, thm:livsic-fullshift}
  Forward: the one-sided measurable headline (\Cref{thm:livsic-measurable-oneside}). Backward:
  vanishing periodic sums feed the substantive backward direction of \Cref{thm:livsic-fullshift} to
  produce a genuine Hölder transfer function, which --- being Hölder --- is continuous, hence
  measurable and an exact (so a.e.) coboundary. Composed with the uniqueness of the a.e.\ solution
  modulo constants (ergodicity of the Bernoulli shift), the measurable solution is a.e.\ the Hölder
  one up to an additive constant
  (\lean{ErgodicTheory.Livsic.measurable_aeCoboundary_ae_eq_holder}).
\end{proof}

\section{The flow tier}

The final layer ports the Livšic obstruction to continuous-time flows. For a one-parameter flow
\(\Phi:\R\to Q\to Q\) and observable \(F:Q\to\R\), the cohomological equation is read integrally.

\begin{definition}[Flow coboundary]\label{def:livsic-flowcob}
  \lean{ErgodicTheory.IsFlowCoboundary}\leanok
  \(F\) is a \emph{flow coboundary} for \(\Phi\) if it is the flow-time derivative (read integrally)
  of some transfer function \(u\):
  \[
    u(\Phi_t q) - u\,q \;=\; \int_0^t F(\Phi_s q)\,ds
    \qquad\text{for all } q,\,t.
  \]
  No integrability is assumed --- the interval integral is read as an opaque real, so the predicate
  is a bare cohomological one mirroring \Cref{def:livsic-coboundary}.
\end{definition}

\begin{theorem}[Flow periodic-orbit obstruction]\label{thm:livsic-flow-obstruction}
  \lean{ErgodicTheory.not_isFlowCoboundary_of_periodicOrbitIntegral_ne_zero}\leanok
  \uses{def:livsic-flowcob}
  If \(\Phi_P q = q\) and the closed-orbit integral \(\int_0^P F(\Phi_s q)\,ds\ne 0\), then \(F\) is
  not a flow coboundary of any transfer function.
\end{theorem}
\begin{proof}\leanok
  \uses{def:livsic-flowcob}
  The fundamental-theorem-of-calculus telescoping around a closed orbit: a transfer function would
  force the orbit integral to equal \(u(\Phi_P q) - u\,q = u\,q - u\,q = 0\). Needs only a bare flow.
\end{proof}

This obstruction is landed on the concrete suspension (mapping-torus) flow \(\zeta_t\) of a base
automorphism \(T\) under a roof \(\tau\), via the cross-section \(x\mapsto[x,0]\).

\begin{definition}[Induced base observable]\label{def:livsic-induced}
  \lean{ErgodicTheory.inducedBaseCocycle}\leanok
  The \emph{induced base observable} of a flow observable \(F\) is the one-lap roof integral
  \(\texttt{inducedBaseCocycle}\ F\ x = \int_0^{\tau x} F([x,s])\,ds\).
\end{definition}

\begin{theorem}[A flow coboundary induces a base coboundary]\label{thm:livsic-induced-cob}
  \lean{ErgodicTheory.inducedBaseCocycle_isCoboundary}\leanok
  \uses{def:livsic-flowcob, def:livsic-induced, def:livsic-coboundary}
  If \(F\) is a flow coboundary of the suspension flow with transfer function \(u\), then its induced
  base observable is a discrete coboundary for the base map \(T\), with transfer function
  \(u\circ(x\mapsto[x,0])\).
\end{theorem}
\begin{proof}\leanok
  \uses{def:livsic-flowcob, def:livsic-induced}
  Evaluate the flow-coboundary equation at the section point \([x,0]\) over one lap \(t = \tau x\):
  flowing that long carries \([x,0]\) to \([Tx,0]\), collapsing the left side to
  \(u([Tx,0]) - u([x,0])\), while the descent \(\zeta_s[x,0] = [x,s]\) identifies the lap integrand
  with the induced observable.
\end{proof}

\begin{theorem}[Tier-1 suspension obstruction]\label{thm:livsic-flow-tier1}
  \lean{ErgodicTheory.not_isFlowCoboundary_of_inducedPeriodicSum_ne_zero}\leanok
  \uses{thm:livsic-induced-cob, thm:livsic-trivial}
  If \(p\) is an \(n\)-periodic base point and the periodic Birkhoff sum of the induced base
  observable does not vanish, then \(F\) is not a flow coboundary of the suspension flow.
\end{theorem}
\begin{proof}\leanok
  \uses{thm:livsic-induced-cob, thm:livsic-trivial}
  Combine \Cref{thm:livsic-induced-cob} with the discrete obstruction certificate of
  \Cref{thm:livsic-trivial} applied to the induced base observable.
\end{proof}

\begin{theorem}[The lap bridge]\label{thm:livsic-lap-bridge}
  \lean{ErgodicTheory.suspension_periodicOrbitIntegral_eq_birkhoffSum}\leanok
  \uses{def:livsic-induced}
  The integral of \(F\) around one full closed flow orbit above an \(n\)-periodic base point equals
  the base Birkhoff sum of the induced base observable, identifying the tier-1 and (flow-native)
  tier-2 obstruction quantities.
\end{theorem}
\begin{proof}\leanok
  \uses{def:livsic-induced}
  The lap decomposition: cut the closed orbit at the successive cross-section returns
  (times \(S_k\tau\,p\)); each lap contributes exactly \(\texttt{inducedBaseCocycle}\ F\ (T^k p)\) by
  a change of variables, and summing the adjacent laps (under per-lap interval integrability)
  reassembles the base Birkhoff sum. The flow-native tier-2 obstruction
  \lean{ErgodicTheory.not_isFlowCoboundary_suspensionFlowMap_of_periodicOrbitIntegral_ne_zero}
  follows by closing a base periodic orbit into a flow-periodic one of period \(S_n\tau\,p\).
\end{proof}

\begin{theorem}[Non-vacuity: the cat-map suspension flow]\label{thm:livsic-cat-suspension}
  \lean{ErgodicTheory.CatMapToral.const_one_not_isFlowCoboundary_catSuspension}\leanok
  \uses{thm:livsic-flow-tier1}
  The constant observable \(1\) is \emph{not} a flow coboundary of the suspension flow over the
  Arnold cat map under the unit roof \(\tau\equiv 1\).
\end{theorem}
\begin{proof}\leanok
  \uses{thm:livsic-flow-tier1}
  Take the fixed point \(0\) of the cat map (base period \(1\)): the induced base observable summed
  around it is the single lap integral \(\int_0^{\tau(0)} 1\,ds = \tau(0) = 1\ne 0\), so the tier-1
  packaged obstruction (\Cref{thm:livsic-flow-tier1}) defeats every transfer function. This exhibits
  a concrete \(F\) with a genuine nonzero closed-orbit integral, certifying the flow obstruction is
  non-vacuous.
\end{proof}

\subsection{The tier-III converse for constant-roof suspension flows}

The obstruction direction above is only half the flow Livšic story. The converse --- vanishing of
all closed-orbit integrals \emph{implies} \(F\) is a flow coboundary --- is delivered here for
\emph{constant-roof} suspension flows (\(\tau\equiv c\)), and delivered \emph{exactly}: because
\texttt{IsFlowCoboundary} is a regularity-free cohomological predicate (\Cref{def:livsic-flowcob}, no
metric on the suspension), the fundamental-domain seam is glued by an \emph{identity}, not a Hölder
estimate. The one input required is that the \emph{base} map discharge its own discrete Livšic
converse for the induced base observable; that is exactly the interface every concrete base in
\Cref{thm:livsic-iff} already supplies.

\begin{lemma}[Seam-glue generator identity]\label{uCover_gen}
  \lean{ErgodicTheory.uCover_gen}\leanok
  \uses{def:livsic-induced}
  Fix a constant roof \(\tau\equiv c\), a base transfer function \(u_0\) cobounding the induced base
  observable, and build the fundamental-domain transfer candidate
  \(\texttt{uCover}(x,s) = u_0(x) + \int_0^{s} F([x,\sigma])\,d\sigma\)
  (\lean{ErgodicTheory.uCover}) lap by lap on the strip. Across one fundamental-domain seam --- the
  identification \([T x, s] = [x, s + c]\) --- the candidate is unchanged, because the base jump
  \(u_0(T x) - u_0(x)\) is \emph{precisely} the lap integral \(\int_0^{c} F([x,\sigma])\,d\sigma\)
  that the fibre integral accumulates over one roof height.
\end{lemma}
\begin{proof}\leanok
  \uses{def:livsic-induced}
  The base cohomological equation \(\texttt{inducedBaseCocycle}\,F\,x = u_0(T x) - u_0(x)\) supplies
  the value of the seam jump; with a constant roof the one-lap integral \(\int_0^{\tau x}F([x,\sigma])
  \,d\sigma = \int_0^{c}F([x,\sigma])\,d\sigma\) is exactly this jump, so the two descriptions of
  \(\texttt{uCover}\) across the seam agree by interval-integral additivity. No estimate; a pure
  identity. Its upgrade to full \(\Z\)-invariance under the suspension action is
  \lean{ErgodicTheory.uCover_act}.
\end{proof}

\begin{theorem}[Tier-III equivalence for constant-roof suspension flows]\label{livsic_suspensionFlow_constRoof}
  \lean{ErgodicTheory.livsic_suspensionFlow_constRoof}\leanok
  \uses{def:livsic-flowcob, def:livsic-induced, thm:livsic-induced-cob, uCover_gen}
  Let \(\Phi = \zeta\) be the suspension flow of \(T\) under a \emph{constant} roof \(\tau\equiv c\),
  let \(F\) be a flow observable with per-fibre interval-integrable restrictions, and suppose the base
  map \(T\) discharges the discrete Livšic converse for the induced base observable --- vanishing
  periodic Birkhoff sums of \(\texttt{inducedBaseCocycle}\,F\) imply it is a discrete coboundary.
  Then
  \[
    F \text{ is a flow coboundary of } \zeta
      \iff \texttt{inducedBaseCocycle}\,F \text{ has vanishing periodic sums.}
  \]
\end{theorem}
\begin{proof}\leanok
  \uses{thm:livsic-induced-cob, uCover_gen}
  Forward is the general obstruction: a flow coboundary induces a base coboundary
  (\Cref{thm:livsic-induced-cob}), which has vanishing periodic sums (\Cref{thm:livsic-trivial}).
  Backward: vanishing periodic sums feed the base converse to produce a base transfer \(u_0\);
  \Cref{uCover_gen} then makes the fundamental-domain candidate seam-invariant, so it descends
  through the suspension quotient (\lean{ErgodicTheory.suspTransfer}, a \(\texttt{Quotient.lift}\))
  to a genuine flow transfer function verifying the continuous-time coboundary equation
  (\lean{ErgodicTheory.suspTransfer_flow}).
\end{proof}

\begin{theorem}[Flow-native form]\label{livsic_suspensionFlow_constRoof_orbitIntegral}
  \lean{ErgodicTheory.livsic_suspensionFlow_constRoof_orbitIntegral}\leanok
  \uses{livsic_suspensionFlow_constRoof, thm:livsic-lap-bridge}
  Under the same hypotheses, the obstruction may be phrased directly in flow-native terms: \(F\) is a
  flow coboundary iff every \emph{closed-orbit integral} of \(F\) vanishes --- for every base
  \(n\)-periodic point \(p\), the integral of \(F\) around the corresponding closed flow orbit (of
  period \(S_n\tau\,p\)) is zero.
\end{theorem}
\begin{proof}\leanok
  \uses{livsic_suspensionFlow_constRoof, thm:livsic-lap-bridge}
  Chain \Cref{livsic_suspensionFlow_constRoof} through the lap-decomposition bridge
  (\Cref{thm:livsic-lap-bridge}), which identifies each closed-orbit flow integral with the base
  Birkhoff sum of the induced observable.
\end{proof}

This equivalence lands on the Arnold cat-map suspension flow, where the discrete base converse is the
already-proved Hölder Livšic theorem for the cat map (\Cref{thm:livsic-cat}).

\begin{theorem}[Flow-Livšic for the cat-map suspension flow]\label{livsic_catSuspensionFlow}
  \lean{ErgodicTheory.CatMapToral.livsic_catSuspensionFlow}\leanok
  \uses{livsic_suspensionFlow_constRoof, thm:livsic-cat}
  For the suspension flow of \(\operatorname{catTorus}\) under the unit roof \(\tau\equiv 1\), a flow
  observable \(F\) whose induced base observable is Hölder (exponent \(0 < r\le 1\)) and whose fibre
  restrictions are interval-integrable is a flow coboundary \emph{iff} every periodic Birkhoff sum of
  its induced base observable vanishes; equivalently (flow-native form
  \lean{ErgodicTheory.CatMapToral.livsic_catSuspensionFlow_orbitIntegral}) iff every closed-orbit
  integral of \(F\) vanishes.
\end{theorem}
\begin{proof}\leanok
  \uses{livsic_suspensionFlow_constRoof, thm:livsic-cat}
  Instantiate \Cref{livsic_suspensionFlow_constRoof} at \(c = 1\), discharging the base converse with
  the Hölder Livšic theorem for the cat map (\Cref{thm:livsic-cat}, \(\texttt{livsic\_catTorus}\)).
\end{proof}

The equivalence has content on \emph{both} sides. The obstruction side is
\Cref{thm:livsic-cat-suspension}: the constant observable \(1\) is not a flow coboundary. The
coboundary side is a concrete zero-mean fibre profile.

\begin{theorem}[Non-vacuity of the coboundary side]\label{isFlowCoboundary_sinFibreObservable}
  \lean{ErgodicTheory.CatMapToral.isFlowCoboundary_sinFibreObservable}\leanok
  \uses{livsic_catSuspensionFlow}
  A \(1\)-periodic fibre profile \(h\) descends to a flow observable \(F[x,s] = h(s)\) on the cat
  suspension (\lean{ErgodicTheory.CatMapToral.fibreObservable}, a \(\texttt{Quotient.lift}\) whose
  well-definedness is the profile's periodicity). For the concrete choice \(h(s) = \sin(2\pi s)\)
  (\lean{ErgodicTheory.CatMapToral.sinFibreObservable}), the induced base observable vanishes
  identically --- its one-lap integral is \(\int_0^{1}\sin(2\pi s)\,ds = 0\) --- so all periodic
  sums vanish and \(\texttt{sinFibreObservable}\) \emph{is} a flow coboundary of the cat-map
  suspension flow.
\end{theorem}
\begin{proof}\leanok
  \uses{livsic_catSuspensionFlow}
  The one-lap integral \(\int_0^{1}\sin(2\pi s)\,ds = 0\) makes \(\texttt{inducedBaseCocycle}\,F\)
  identically zero, so its periodic Birkhoff sums vanish trivially; the backward direction of
  \Cref{livsic_catSuspensionFlow} supplies the flow transfer function. Together with the obstruction
  witness \(F = 1\) (\Cref{thm:livsic-cat-suspension}), this certifies the tier-III equivalence is
  non-vacuous on both sides --- some observables cobound, others do not.
\end{proof}

\paragraph{The remaining wall.} The constant-roof converse above is exact; what is \emph{not}
delivered is the \emph{variable-roof} flow Livšic theorem, where the fundamental-domain seam
\(s = \tau(x)\) moves with the base point and the one-lap identity of \Cref{uCover_gen} no longer
glues the transfer function by a bare cohomological equation --- a genuine Hölder-regularity gluing
would be required. It is recorded here as the disclosed open front remaining after the constant-roof
equivalence \Cref{livsic_suspensionFlow_constRoof}.
